\begin{section}{The Real Field}

There exists an ordered field $(\R, +, \cdot)$ which has the
\emph{least upper bound property}. Moreover, $\R$ contains $\Q$ as a subfield, i.e.,
$\Q \subseteq \R$ and the operations of addition and multiplication on $\R$ coincide with those on $\Q$.
\\[6mm]
We can construct $\R$ from $\Q$. The construction of $\R$ is rather complicated, so it will be discussed
separately in \hyperref[chap:ConstructionOfR]{Appendix}. In this section, we will just assume that $\R$ exists.

\bigskip

\begin{theorem} \label{thm:ArchimedeanProperty}
    For any $x\in \R^{+}$ and $y \in \R$, there exists some $n \in \N$ such that $nx > y$.
    This theorem is called the \textbf{Archimedean property} of $\R$.
\end{theorem}

\begin{proof}
    Let $x \in \R^{+}$ and $y \in \R$ be given. Let $E$ be defined as follows:
    \[
    E \coloneqq \setbuilder*{nx}{nx \leq y, \; n \in \N}.
    \]
    If $E = \vnot$, which implies that $x > y$, then $n = 1$ satisfies $nx > y$. Otherwise, $E$ is nonempty and bounded
    above by $y$. By the least upper bound property, there exists $\alpha = \sup E \in \R$. Then we have the following:
    \begin{itemize}
        \item $\alpha - x$ is not an upper bound of $E \implies$ There exists some $m \in \N$ such that $mx > \alpha - x$.
        \item $mx > \alpha - x \implies (m + 1)x > \alpha \implies (m + 1)x \notin E \implies (m + 1)x > y$.
    \end{itemize}
    Therefore, we can take $n = m + 1$ to complete the proof.
\end{proof}

\bigskip

\begin{corollary} \label{cor:ArchimedeanPropertyCorollary}
    For any $\veps \in \R^{+}$ and $b \in (1, + \infty)$, there exists some $n \in \N$ such that $b^{-n} < \veps$.
\end{corollary}

\begin{proof}
    For any $n \in \N$ and $x \in [-1, + \infty)$, we have
    \[
    (1 + x)^{n} \geq 1 + nx.
    \]
    Let $a \in \R^{+}$ and $\veps \in \R^{+}$ be given. By \Cref{thm:ArchimedeanProperty},
    there exists some $n \in \N$ such that $na > \dfrac{1}{\veps}$. Then we have
    \[
    (1 + a)^{n} \geq 1 + na > na > \frac{1}{\veps} \implies (1 + a)^{-n} < \veps.
    \]
    Since $b > 1$, we can take $a = b - 1 > 0$ to complete the proof.
\end{proof}

\bigskip

\begin{theorem} \label{thm:BoundedSubsetOfZInRHasExtremum}
    Any nonempty subset of $\Z$ which is bounded below in $\R$ has a unique least element,
    and any nonempty subset of $\Z$ which is bounded above in $\R$ has a unique greatest element.
\end{theorem}

\begin{proof}
    Let $E$ be a nonempty subset of $\Z$ which is bounded below in $\R$. Then there exists a lower bound $\alpha \in \R$
    of $E$. Then we have the following:
    \begin{itemize}
        \item If $\alpha \geq 0$, then $-1 \in \Z$ is a lower bound of $E$.
        \item If $\alpha < 0$, then by \Cref{thm:ArchimedeanProperty}, there exists some $n \in \N$
        such that $n > - \alpha$, which implies that $-n \in \Z$ is a lower bound of $E$.
    \end{itemize}
    Therefore, $E$ is bounded below in $\Z$. By \Cref{thm:BoundedSubsetOfZinZHasExtremum}, $E$ has a unique least element.
    \\[6mm]
    Let $F$ be a nonempty subset of $\Z$ which is bounded above in $\R$. Then there exists an upper bound $\beta \in \R$
    of $F$. Then we have the following:
    \begin{itemize}
        \item By \Cref{thm:ArchimedeanProperty}, there exists some $n \in \N$
        such that $n > \beta$, which implies that $n \in \Z$ is an upper bound of $F$.
    \end{itemize}
    Therefore, $F$ is bounded above in $\Z$. By \Cref{thm:BoundedSubsetOfZinZHasExtremum}, $F$ has a unique greatest element.
\end{proof}

\bigskip

\begin{definition} \label{def:FloorCeilingFunctions}
    For any $x \in \R$, the \textbf{floor function} $\floor{x}$
    and \textbf{ceiling function} $\ceil{x}$ are defined as follows:
    \begin{itemize}
        \item $\floor{x} \coloneqq \max\setbuilder*{k \in \Z}{k \leq x}$.
        \item $\ceil{x} \coloneqq \min\setbuilder*{k \in \Z}{k \geq x}$.
    \end{itemize}
    \Cref{thm:BoundedSubsetOfZInRHasExtremum} guarantees the
    \emph{existence} and \emph{uniqueness} of $\floor{x}$ and $\ceil{x}$.
\end{definition}

\bigskip

\begin{theorem} \label{thm:QIsDenseInR}
    For any $- \infty < x < y < + \infty$, there exists some $p \in \Q$ such that $x < p < y$.
    This theorem states that $\Q$ is \textbf{dense} in $\R$.
\end{theorem}

\begin{proof}
    Let $\veps \coloneqq y - x > 0$. By \Cref{thm:ArchimedeanProperty}, there exists some $n \in \N$
    such that $\dfrac{1}{n} < \veps$. Let $S$ be defined as follows:
    \[
    S \coloneqq \setbuilder*{k \in \Z}{\dfrac{k}{n} < x}.
    \]
    $S$ is nonempty since $n\floor{x} - 1 \in S$, and is bounded above by $nx \in \R$.
    By \Cref{thm:BoundedSubsetOfZInRHasExtremum}, $S$ has a greatest element, say $m$.
    \[
    x < \frac{m + 1}{n} < x + \frac{1}{n} < x + (y - x) = y.
    \]
    Therefore, we can take $p = \dfrac{m + 1}{n} \in \Q$ to complete the proof.
\end{proof}

\bigskip

\begin{theorem} \label{thm:ExistenceOfNthRoots}
    For each $x \in \R^{+}$ and $n \in \N$, there exists a unique $y \in \R^{+}$ such that
    $y^{n} = x$. We denote $y$ by $\sqrt[n]{x}$ or $x^{\frac{1}{n}}$.
\end{theorem}

\begin{proof}
    Let $E$ be defined as follows:
    \[
    E \coloneqq \setbuilder*{t \in \R^{+}}{t^{n} < x}.
    \]
    $E$ is nonempty since $2^{-1} \in E$ if $x\geq 1$, and $x^{2} \in E$ otherwise.
    Also, $E$ is bounded above by $\max\set*{x, 1}$. By the least upper bound property,
    there exists $\alpha = \sup E \in \R^{+}$. We will show that $\alpha^{n} = x$.
    \\[6mm]
    For any $a, b \in \R^{+}$ with $a < b$, we have
    \[
    b^{n} - a^{n} = (b - a)\sum_{k = 0}^{n - 1} b^{n - 1 - k} a^{k} < (b - a)nb^{n - 1}.
    \]
    Assume that $\alpha^{n} < x$. By \Cref{thm:QIsDenseInR}, there exists some $h \in \R^{+}$ such that
    \[
    h < \min\set*{1, \; \frac{x - \alpha^{n}}{n(\alpha + 1)^{n - 1}}}.
    \]
    Put $a = \alpha$ and $b = \alpha + h$. Then we have
    \begin{align*}
        &(\alpha + h)^{n} - \alpha^{n} < hn(\alpha + h)^{n - 1} < hn(\alpha + 1)^{n - 1} \\
        \implies &(\alpha + h)^{n} < \alpha^{n} + hn(\alpha + 1)^{n - 1} < \alpha^{n} + (x - \alpha^{n}) = x.
    \end{align*}
    Therefore, we have $\alpha + h \in E$, which is a contradiction.
    \\[6mm]
    Assume that $\alpha^{n} > x$. By \Cref{thm:QIsDenseInR}, there exists some $h \in \R^{+}$ such that
    \[
    h < \min\set*{1, \; \alpha, \; \frac{\alpha^{n} - x}{n(\alpha + 1)^{n - 1}}}.
    \]
    Put $a = \alpha - h$ and $b = \alpha$. Then we have
    \begin{align*}
        &\alpha^{n} - (\alpha - h)^{n} < hn\alpha^{n - 1} < hn(\alpha + 1)^{n - 1} \\
        \implies &(\alpha - h)^{n} > \alpha^{n} - hn(\alpha + 1)^{n - 1} > \alpha^{n} - (\alpha^{n} - x) = x.
    \end{align*}
    Therefore, we have $\alpha - h$ is an upper bound of $E$, which is a contradiction.
    \\[6mm]
    Therefore, we have $\alpha^{n} = x$, which guarantees the \emph{existence} of $y$. For any $\alpha, \beta \in \R^{+}$,
    we have $\alpha \neq \beta \implies \alpha^{n} \neq \beta^{n}$, which guarantees the \emph{uniqueness} of $y$.
\end{proof}

\bigskip

\begin{definition} \label{def:GreedyBaseExpansionOfRealNumbers}
    Let $y \in [0, 1)$ and $b \in \N$ be given with $b \geq 2$. A sequence $\paren*{a_{n}}_{n = 1}^{\infty}$
    is called the \textbf{greedy base $b$ expansion} of $y$ if for each $n \in \N$, we have
    \[
    0 \leq y - \sum_{k = 1}^{n} \frac{a_{k}}{b^{k}} < \frac{1}{b^{n}}
    \]
    where $a_{k} \in S_{b} \coloneqq \set*{0, 1, 2, \dots, b - 1}$ for all $k \in \N$.
    Especially, if $b = 10$, the sequence is called the \textbf{greedy decimal expansion} of $y$.
\end{definition}

\bigskip

\begin{theorem} \label{thm:ExistenceAndUniquenessOfGreedyBaseExpansion}
    For any $y \in [0, 1)$, the greedy base $b$ expansion of $y$ exists and is unique. 
\end{theorem}

\begin{proof}
    Let $\paren*{t_{n}}_{n = 1}^{\infty}$, $\paren*{T_{n}}_{n = 1}^{\infty}$
    and $\paren*{a_{n}}_{n = 1}^{\infty}$ be defined as follows:
    \begin{itemize}
        \item $t_{1} \coloneqq by$, $T_{1} \coloneqq \setbuilder*{t \in S_{b}}{t \leq t_{1}}$, $a_{1} \coloneqq \max T_{1}$.
        \item $t_{2} \coloneqq b^{2}y - ba_{1}$, $T_{2} \coloneqq \setbuilder*{t \in S_{b}}{t \leq t_{2}}$,
        $a_{2} \coloneqq \max T_{2}$.
        \item[] $\vdots$ 
        \item $t_{n} \coloneqq b^{n}y - \displaystyle \sum_{k = 1}^{n - 1} a_{k}b^{n - k}$,
        $T_{n} \coloneqq \setbuilder*{t \in S_{b}}{t \leq t_{n}}$, $a_{n} \coloneqq \max T_{n}$
        
        for all $n \in \N$ with $n \geq 2$,
    \end{itemize}
    Assume that $0 \leq t_{n} < b$ for some $n \in \N$. Then, since $0 \in T_{n}$ and $T_{n}$ is bounded above, $T_{n}$ has
    a greatest element by \Cref{thm:BoundedSubsetOfZinZHasExtremum}, which implies that $\paren*{a_{n}}_{n = 1}^{\infty}$
    is well-defined. Let $P(n)$ be the statement that
    \[
    P(n) : 0 \leq t_{n} < b \text{ and } a_{n} \text{ is well-defined}.
    \]
    \begin{itemize}
        \item \textbf{Base case:} $P(1)$ is true since $0 \leq t_{1} = by < b$.
        \item \textbf{Inductive step:} Assume that $P(n)$ is true for some $n \in \N$. Then we have
        \begin{align*}
            t_{n + 1} &= b^{n + 1}y - \sum_{k = 1}^{n} a_{k}b^{n + 1 - k} \\
            &= b\paren*{b^{n}y - \sum_{k = 1}^{n} a_{k}b^{n - k}} \\
            &= b\paren*{b^{n}y - \sum_{k = 1}^{n - 1} a_{k}b^{n - k} - a_{n}} \\
            &= b\paren*{t_{n} - a_{n}}.
        \end{align*}
        Since $a_{n} = \max T_{n}$ and $t_{n} < b$, we have
        \[
        a_{n} \leq t_{n} < a_{n} + 1 \implies 0 \leq t_{n} - a_{n} < 1 \implies 0 \leq t_{n + 1} = b(t_{n} - a_{n}) < b.
        \]
        Thus $P(n + 1)$ is true whenever $P(n)$ is true.
    \end{itemize}
    By the \emph{principle of mathematical induction}, $P(n)$ is true for all $n \in \N$.